\begin{helper}
Let $\mathcal F$ be a finite $k$-uniform multi-hypergraph, let $f$ be an
edge, and suppose every vertex of \(f\) has degree exactly \(D\).  If an
edge \(g\ne f\) contains two distinct vertices \(u,v\in f\), then the number
of line-graph neighbors of \(f\) is at most
\[
  k(D-1)-1 .
\]
\end{helper}
\begin{proof}
By \noderef{LineGraphOfHypergraph}, the line-graph neighbors of \(f\) are
exactly the edge labels \(e\ne f\) whose hyperedges meet \(f\).  Let this
finite set be \(N\).

For each vertex \(x\in f\), the number of edge labels \(e\ne f\) containing
\(x\) is \(D-1\): by the saturation hypothesis and
\noderef{HypergraphDegree}, exactly \(D\) edge labels contain \(x\), and one
of them is \(f\).  Since \(\mathcal F\) is \(k\)-uniform by
\noderef{UniformHypergraph}, the total number of incidences
\[
  (x,e)\quad\text{with }x\in f,\ e\ne f,\ x\in e
\]
is therefore \(k(D-1)\).

The same incidence set can be counted by summing, over \(e\in N\), the size of
\(f\cap e\).  Every \(e\in N\) contributes at least one incidence, because it
meets \(f\).  The distinguished edge \(g\) contributes at least two
incidences, namely the two distinct vertices \(u\) and \(v\).  Hence the
incidence count is at least \(|N|+1\).  Combining this lower bound with the
previous exact count gives
\[
  |N|+1 \le k(D-1),
\]
and therefore \(|N|\le k(D-1)-1\).  This is precisely the asserted upper
bound on the line-graph neighbor count of \(f\).
\end{proof}
